\documentclass[10pt,a4paper]{amsart}
\usepackage[margin=19mm]{geometry}
\usepackage{amsmath,amssymb,mathtools}
\usepackage[scr=rsfs,cal=euler]{mathalfa}
\usepackage{xfrac}
\usepackage[unicode,hidelinks,bookmarks=false]{hyperref}
\newcommand{\ie}{i.e.\ }
\newcommand{\cf}{cf.\ }
\newcommand{\resp}{respectively\ }
\newcommand{\Subsection}{\subsection}
\providecommand{\nobreakdash}{\nobreak\hbox{-}}
\setlength{\emergencystretch}{2em}
\multlinegap0pt
\usepackage{enumitem}
\usepackage{amscd}
\DeclareMathOperator{\dist}{dist}
\DeclareMathOperator{\intr}{int}
\DeclareMathOperator{\conv}{co}
\DeclareMathOperator{\intt}{int}
\DeclareMathOperator{\LL}{\mathcal{L}}
\DeclareMathOperator{\argmin}{argmin}
\DeclareMathOperator{\dom}{dom}
\DeclareMathOperator{\Min}{Min}
\DeclareMathOperator{\cl}{cl}
\DeclareMathOperator{\bd}{bd}
\DeclareMathOperator{\co}{co}
\DeclareMathOperator{\lin}{lin}
\DeclareMathOperator{\aff}{aff}
\DeclareMathOperator{\ri}{ri}
\DeclareMathOperator{\rint}{rint}
\DeclareMathOperator{\rext}{rext}
\DeclareMathOperator{\diam}{diam}
\DeclareMathOperator{\gph}{gph}
\DeclareMathOperator{\WMin}{WMin}
\DeclareMathOperator{\reg}{reg}
      \setcounter{page}{1}
\newtheorem{definition}{Definition}[section]
\newtheorem{theorem}{Theorem}[section]
\newtheorem{corollary}{Corollary}[section]
\newtheorem{lemma}{Lemma}[section]
\newtheorem{proposition}{Proposition}[section]
\newtheorem{remark}{Remark}[section]
\newtheorem{example}{Example}[section]
\newtheorem{counterexample}{Counterexample}[section]
\renewcommand{\baselinestretch}{1.25}
\begin{document}
\title[Directional Derivatives and Error Bounds of Merit Functions ]{Directional Derivatives and Error Bounds of Merit Functions in Vector Optimization}
\author{Nanchang, Jiangxi 330013, China; Yu Han}
\date{}
\maketitle
\noindent\emph{Source and licence.} Directional Derivatives and Error Bounds of Merit Functions in Vector Optimization. Version arXiv:2607.18781v1 (CC BY 4.0). This material is distributed under the Creative Commons Attribution 4.0 International licence, http://creativecommons.org/licenses/by/4.0/, as recorded in the arXiv OAI licence metadata for this exact version and shown on the abstract page https://arxiv.org/abs/2607.18781v1. Full rights notice: licenses/arxiv-2607.18781-CC-BY-4.0.txt.\par\smallskip
\noindent\textit{Editorial scope.} Counted unit: the original \texttt{theorem} environment \texttt{XLWCJK} at source lines 1025--1093 together with its complete proof at lines 1095--1569; the delivered document keeps source lines 154--1570 so that every referenced label and every citation resolves inside it. Native editable source is the paper's own concatenated TeX; the AMS wrapper is editorial formatting only. No publisher class, upstream script or unrelated artwork is redistributed.
\begin{equation}\label{eq:VP}
	\operatorname{Min}_C \; F(x) \quad \text{subject to} \quad x \in A,
\end{equation}
where $\operatorname{Min}_C$ denotes minimization with respect to the partial order $\leq_C$ induced by $C$, i.e., $y_1 \leq_C y_2$ if and only if $y_2 - y_1 \in C$. The strict partial order $\ll_C$ is defined by $y_1 \ll_C y_2$ if and only if $y_2 - y_1 \in \intr C$.

A feasible point $\bar{x} \in A$ is called a \emph{weakly efficient solution} of \eqref{eq:VP} if there exists no $x \in A$ such that $F(x) \ll_C F(\bar{x})$, which is equivalent to
\[
F(\bar{x}) \notin F(A) + \intr C.
\]
The set of all weakly efficient solutions is denoted by $E_w(\text{VP})$ or simply $E_w$. Its image $F(E_w)$ consists of the \emph{weakly efficient points} of $F(A)$, denoted by $\operatorname{WMin}(F(A), C)$.  


For a subset $A \subseteq X$ with respect to a closed convex cone $P \subseteq X$ ($\intr P \neq \emptyset$), the set of \emph{weakly efficient points} of $A$ is
\[
\operatorname{WMin}(A, P) := \{a \in A : a \notin A + \intr P\}.
\]




A set $A$ is said to have the  weak domination property  (WDP) with respect to $P$, if for every $z \in A \setminus \operatorname{WMin}(A, P)$, there exists $z' \in \operatorname{WMin}(A, P)$ such that $z' \leq_P z$.

\begin{remark} \label{RngXEQ}     It is easy to verify that  $A$ has  the  weak domination property with respect to $P$ if and only if for any $z \in A \setminus \operatorname{WMin}(A, P)$, there exists $z' \in \operatorname{WMin}(A, P)$ such that $z'  \ll_P  z$.
\end{remark}


The canonical bilinear pairing is denoted by $\langle \cdot, \cdot \rangle$: for $x^* \in X^*$ and $x \in X$, $\langle x^*, x \rangle := x^*(x)$. The norm on $X$ (and other spaces) is denoted by $\|\cdot\|$. The closed unit ball and unit sphere of $X$ are
\[
B_X := \{x \in X : \|x\| \leq 1\}, \quad S_X := \{x \in X : \|x\| = 1\}.
\]

$B(x_0,r)$ denotes the open ball with center $x_0$ and radius $r > 0$. For a subset $D \subseteq X$, its interior, boundary, closure, convex hull  and  closed convex hull    are denoted by $\intr D$, $\bd D$, $\cl D$, $\co D$ and $\overline{\co} D$,  respectively.

 


For a nonempty set \(D \subseteq X\), the distance from a point \(x \in X\) to the set \(D\) is defined as
\[
d_D(x) := \inf_{y \in D} \|x - y\|,
\]
with the convention that \(d_\emptyset(\cdot) = +\infty\). Occasionally, we use \(d(x,D)\) to denote the distance from a point \(x \in X\) to the set \(D\).
The distance between two nonempty sets \(Q \subseteq X\) and \(D \subseteq X\) is defined by
\[
d(Q,D) = \inf \{ \|x - y\| : x \in Q,\; y \in D \}.
\]
The Hausdorff distance between $Q$ and $D$ is defined by
$$d_H \left( {Q,D} \right): = \max \left\{ {g\left( {Q,D} \right),g\left( {D,Q} \right)} \right\},$$
where $g\left( {Q,D} \right): = \mathop {\sup }\limits_{x \in Q} d\left( {x,D} \right)$.  
 





The  indicator function  of a  nonempty set   $D$ is
\[
\delta_D(x) := \begin{cases}
	0, & x \in D, \\
	+\infty, & x \notin D.
\end{cases}
\]


For  $ x^* \in X^*$,     the support function   of a  nonempty set   $D$ is
$$\sigma_D(x^*)=\sup_{y\in D}  \;  \langle x^*,y\rangle.$$

The \emph{projection} of a point $x$ onto a closed convex set $D$ is the (possibly empty) set
\[
P_D(x) := \{\bar{x} \in D : \|x - \bar{x}\| = d_D(x)\}.
\]
When $X$ is a Hilbert space, $P_D(x)$ is a singleton for every $x \in X$.

 

The \emph{positive polar cone} (or dual cone) of a set $Q \subseteq X$ is
\[
Q^+ := \{x^* \in X^* : \langle x^*, x \rangle \geq 0, \; \forall x \in Q\},
\]
and the \emph{negative polar cone} is $Q^\ominus := -Q^+ = \{x^* \in X^* : \langle x^*, x \rangle \leq 0, \; \forall x \in Q\}$.    The unit sphere of \(Q^+\) is denoted by \(S(Q^+)\), i.e., \(S(Q^+) = \{ x^* \in Q^+ : \|x^*\| = 1 \}\).

 
\begin{definition}  
 For a nonempty set $D \subseteq X$ and $x \in D$, the \emph{tangent cone} and \emph{normal cone} of  $D$ at  $x$ are defined, respectively, by  
 \begin{align*}
 	T_D(x) &:= \bigl\{ d \in X : \exists\, t_k \downarrow 0,\; \exists\, d_k \to d \text{ with } x + t_k d_k \in D \bigr\},\\
 	N_D(x) &:= \{x^* \in X^* : \langle x^*, h \rangle \leq 0, \; \forall h \in T_D(x)\}.
 \end{align*}
  \end{definition}

 

\begin{remark} \label{RtDFxl}    \cite{Zalinescu2002, AE1984}    Let  $D \subseteq X$ be a closed convex set  and $x \in D$.  Then
	\begin{align*}
		T_D(x) &= \cl \left( \bigcup_{t > 0} t(D - x) \right),\\
		N_D(x) &= \{ x^* \in X^* : \langle x^*, y - x \rangle \le 0,\ \forall y \in D \}.  
	\end{align*}
	Moreover, we have  $N_D(x) = T_D(x)^\ominus$ and $T_D(x) = N_D(x)^\ominus$.
\end{remark}


 

\subsection{The  oriented distance function}

The oriented distance function (see \cite{Hiriart1,Hiriart2}) of a set $D \subseteq Y$ is defined as
$$
\Delta_D(y) := d_D(y) - d_{Y \setminus D}(y) = \begin{cases}
	d_D(y), & y \in Y \setminus D, \\
	-d_{Y \setminus D}(y), & y \in D.
\end{cases}
$$
With the convention: if $D = \emptyset$, then $\Delta_D \equiv +\infty$; if $D = Y$, then $\Delta_D \equiv -\infty$.

Next,we collect the basic properties of the oriented distance function.
\begin{proposition}  (see \cite{XuLi2016, KTZ2015, Zaffaroni2003})  \label{XSQ}
	Let $D \subseteq Y$ be nonempty and $D \neq Y$. Then:
	\begin{itemize}
		\item[(i)] $\Delta_D$ is real-valued and $\Delta_{Y \setminus D} = -\Delta_D$;
		\item[(ii)] $\Delta_D$ is $1$-Lipschitz;
		\item[(iii)] $\Delta_D(y) < 0 \Leftrightarrow y \in \operatorname{int} D$;
		\item[(iv)] $\Delta_D(y) = 0 \Leftrightarrow y \in \partial D$;
		\item[(v)] $\Delta_D(y) > 0 \Leftrightarrow y \in \operatorname{int} D^c$;
		\item[(vi)] If $D$ is closed, then $D = \{y \in Y : \Delta_D(y) \leq 0\}$;
		\item[(vii)] If $D$ is a cone, then $\Delta_D$ is positively homogeneous;
		\item[(viii)] If $D$ is convex, then $\Delta_D$ is a convex function;
		\item[(ix)] If $D$ is a convex cone, then $\Delta_D$ is nonincreasing with respect to the order induced by $D$ on $Y$, i.e., for any $y_1, y_2 \in Y$,
		\[
		y_1 - y_2 \in D \;\Longrightarrow\; \Delta_D(y_1) \leq \Delta_D(y_2);
		\]
		if $D$ has nonempty interior, then
		\[
		y_1 - y_2 \in \operatorname{int} D \;\Longrightarrow\; \Delta_D(y_1) < \Delta_D(y_2).
		\]
	\end{itemize}
\end{proposition}




By Lemma 2 (ii) of \cite{LiuNgYang2009}, we have the following lemma.

\begin{lemma}  \label{QWED} Assume that  $D \subseteq Y$  is nonempty and convex and  $y \in Y$. If ${d_D}\left( y \right) > 0$,  then
	$${d_D}\left( y \right) = \mathop {\sup }\limits_{{y^*} \in {S_{{Y^*}}}} \mathop {\inf }\limits_{u \in D} \left\langle {{y^*},y - u} \right\rangle .$$	
\end{lemma}

By Lemma  \ref{QWED} and the proof of Proposition 1  of \cite{LiuNgYang2009}, we   get the following lemma.  
\begin{lemma}  \label{QWAD} Assume that  $D \subseteq Y$  is nonempty, closed and convex, $D \ne Y$
	and  $y \in D$. Then
	$${d_{Y\backslash D}}(y) = \mathop {\inf }\limits_{{y^*} \in {S_{{Y^*}}}} \mathop {\sup }\limits_{u \in D} \left\langle {{y^*},y - u} \right\rangle .$$
\end{lemma}
 

\begin{proposition} \label{wqaz} Let $D$ be a nonempty and convex subset of $Y$. Then 
	$${\Delta _D}(y) = \mathop {\sup }\limits_{{y^*} \in {S_{{Y^*}}}} \mathop {\inf }\limits_{u \in D} \left\langle {{y^*},y - u} \right\rangle , \quad \forall y \in Y.$$	
\end{proposition} 
\begin{proof}  It follows from      Lemma 2.3 of \cite{Han2022} that ${\Delta _{{\rm{cl}}D}}(y) = {\Delta _D}(y)$.  There are two cases to be considered.

Case 1. $y \notin {\rm{cl}}D$. Then we have ${d_D}\left( y \right) > 0$ and ${\Delta _D}\left( y \right) = {d_D}\left( y \right)$. This together with  Lemma  \ref{QWED}  implies that 
$${\Delta _D}\left( y \right) = {d_D}\left( y \right) = \mathop {\sup }\limits_{{y^*} \in {S_{{Y^*}}}} \mathop {\inf }\limits_{u \in D} \left\langle {{y^*},y - u} \right\rangle .$$	

Case 2. $y \in {\rm{cl}}D$.  This yields    that ${\Delta _{{\rm{cl}}D}}(y) =  - {d_{Y\backslash {\rm{cl}}D}}(y)$. By Lemma  \ref{QWAD}, we have 
$${d_{Y\backslash {\rm{cl}}D}}(y) = \mathop {\inf }\limits_{{y^*} \in {S_{{Y^*}}}} \mathop {\sup }\limits_{u \in {\rm{cl}}D} \left\langle {{y^*},y - u} \right\rangle = \mathop {\inf }\limits_{{y^*} \in {S_{{Y^*}}}} \mathop {\sup }\limits_{u \in D} \left\langle {{y^*},y - u} \right\rangle .$$
Combining this with  $ - {S_{{Y^*}}} = {S_{{Y^*}}}$, we get
\begin{eqnarray*} {\Delta _D}(y) &=& {\Delta _{{\rm{cl}}D}}(y) =  - {d_{Y\backslash {\rm{cl}}D}}(y) = - \mathop {\inf }\limits_{{y^*} \in {S_{{Y^*}}}} \mathop {\sup }\limits_{u \in D} \left\langle {{y^*},y - u} \right\rangle \\
	&=& \mathop {\sup }\limits_{{y^*} \in {S_{{Y^*}}}} \mathop {\inf }\limits_{u \in D} \left\langle { - {y^*},y - u} \right\rangle = \mathop {\sup }\limits_{{y^*} \in  - {S_{{Y^*}}}} \mathop {\inf }\limits_{u \in D} \left\langle {{y^*},y - u} \right\rangle  \\
	&=& \mathop {\sup }\limits_{{y^*} \in {S_{{Y^*}}}} \mathop {\inf }\limits_{u \in D} \left\langle {{y^*},y - u} \right\rangle .		
\end{eqnarray*}
 \end{proof}  


\begin{remark}  In view of Proposition \ref{wqaz},       the condition \(\operatorname{ri}(A) \neq \emptyset\) in Proposition 1 of \cite{LiuNgYang2009} is redundant,  the conclusion remains valid even if this condition is removed.
   \end{remark}

By Proposition \ref{wqaz}, the following corollary is readily obtained.
 \begin{corollary}\label{cor:Delta}
 	For a nonempty and  convex cone $P \subseteq Y$,  we have 
 $$
 		\Delta_P(y) = \sup_{y^* \in S(P^+)} \langle -y^*, y \rangle, \quad \forall y \in Y.
$$
 \end{corollary}


\begin{remark}   As a consequence of Corollary \ref{cor:Delta}, the hypothesis $\intr K \neq \emptyset$ in Corollary 2 of \cite{LiuNgYang2009} turns out to be redundant, the conclusion remains valid even without this condition.
  \end{remark}


 
 

\subsection{Problem-specific notation}

Returning to the vector optimization problem \eqref{eq:VP}, define the lifted cone
\[
C_X := F^{-1}(C) = \{x \in X : F(x) \in C\}.
\]
One readily verifies that if \( F(X) \cap \operatorname{int} C \neq \emptyset \), then \( \operatorname{int} C_X = F^{-1}(\operatorname{int} C) \neq \emptyset \).  
To avoid the triviality we always assume that $F(X) \cap \intr C \neq \emptyset$.



The adjoint operator $F^*: Y^* \to X^*$ is defined by $\langle F^*(y^*), x \rangle = \langle y^*, F(x) \rangle$, and satisfies the key relation (see \cite{Zalinescu1978}):
$$
	F^*(C^+) = (F^{-1}(C))^+ = C_X^+.
$$

The set
$$
	\widehat{A} := X \setminus (A + \intr C_X)
$$
will play a central role. Note that $E_w  = A \cap \widehat{A}$ (see \cite[Theorem 1]{LiuNgYang2009}).
It is clear that  $E_w $ is closed.

Let $K := \overline{\operatorname{co}}^{w^*}\big(S(C^+)\big)$,   where $S(C^+)=\{y^*\in C^+:\|y^*\|=1\}$  and ${\overline{\operatorname{co}}}^{w^*}$ denotes the closed convex hull in the weak$^*$ topology.
   In this paper, the set $K$ is of central importance, as many of the results are formulated in terms of $K$.




\begin{proposition}   \label{RDCxl}   
	Let $K = \overline{\operatorname{co}}^{w^*}\big(S(C^+)\big)$. Then $K$ is a nonempty weak$^*$-compact convex set  and  $0 \notin K$ 
   \end{proposition}
\begin{proof}  
Clearly, $K \subseteq C^+$ and  $K \subseteq \{y^*\in Y^*:\|y^*\|  \leq 1\}$.
Since $S(C^+)=\{y^*\in C^+:\|y^*\|=1\}$ is nonempty weak$^*$-compact, we obtain that $K$ is a nonempty weak$^*$-compact convex set.  

Fix \(y_0 \in \operatorname{int} C\). Then  there exists \(\varepsilon > 0\) such that
\begin{equation}\label{Edfty1}
	y_0 + \varepsilon B_Y \subseteq C,
\end{equation}
where \(B_Y \) is the closed unit ball in \(Y\).
Take an arbitrary \(y^* \in S(C^+)\), i.e., \(y^* \in C^+\) and \(\|y^*\| = 1\). By the definition of \(C^+\), we have \(\langle y^*, c \rangle \ge 0\) for all \(c \in C\).
In particular, for every \(u \in B_Y\), by~\eqref{Edfty1} we have \(y_0 + \varepsilon u \in C\). Hence
\[
\langle y^*, y_0 + \varepsilon u \rangle \ge 0, \quad \forall u \in B_Y.
\]
This is equivalent to
\[
\langle y^*, y_0 \rangle \ge -\varepsilon \langle y^*, u \rangle \quad \forall u \in B_Y.
\]
Taking the supremum over the right-hand side, we obtain
\begin{equation}\label{Edfty2}
	\langle y^*, y_0 \rangle \ge \sup_{u \in B_Y} \bigl(-\varepsilon \langle y^*, u \rangle\bigr)
	= \varepsilon \sup_{u \in B_Y} \bigl(-\langle y^*, u \rangle\bigr).
\end{equation}
Note that
\[
\sup_{u \in B_Y} \bigl(-\langle y^*, u \rangle\bigr) = \sup_{u \in B_Y} \langle y^*, -u \rangle = \sup_{v \in B_Y} \langle y^*, v \rangle = \|y^*\| = 1,
\]
which together with~\eqref{Edfty2} yields
\[
\langle y^*, y_0 \rangle \ge \varepsilon \cdot 1 = \varepsilon > 0.
\]
Since \(y^* \in S(C^+)\) is arbitrary, we conclude that
\begin{equation}\label{Edfty3}
	\inf_{y^* \in S(C^+)} \langle y^*, y_0 \rangle \ge \varepsilon > 0.
\end{equation}

Now take any \(v \in \operatorname{co}(S(C^+))\). Then there exist finitely many scalars \(\lambda_i \ge 0\) with \(\sum_{i=1}^n \lambda_i = 1\) and \(y_i^* \in S(C^+)\) such that
\[
v = \sum_{i=1}^n \lambda_i y_i^*.
\]
For each \(i\), by~\eqref{Edfty3} we have \(\langle y_i^*, y_0 \rangle \ge \varepsilon\). Therefore,
\begin{equation}\label{Edfty4}
	\langle v, y_0 \rangle = \sum_{i=1}^n \lambda_i \langle y_i^*, y_0 \rangle \ge \sum_{i=1}^n \lambda_i \varepsilon = \varepsilon.
\end{equation}

If \(w = 0 \in K\), there exists a net \(\{v_\alpha\} \subseteq \operatorname{co}(S(C^+))\) such that \(v_\alpha \xrightarrow{w^*} w\). Then \(\langle v_\alpha, y_0 \rangle \to \langle w, y_0 \rangle\).

By~\eqref{Edfty4}, for every \(\alpha\) we have \(\langle v_\alpha, y_0 \rangle \ge \varepsilon\). Hence,
\[
0 = \langle 0, y_0 \rangle = \langle w, y_0 \rangle = \lim_\alpha \langle v_\alpha, y_0 \rangle \ge \varepsilon > 0,
\]
which is a contradiction. Therefore, \(0 \notin K\).
   \end{proof}  

\section{Properties and dual representation of the merit function  }\label{sec:merit}

 The aim of this section is to establish the fundamental analytic properties of 
 $\theta$—its concavity, Lipschitz continuity, and dual representation—which will be indispensable for the subsequent study of directional derivatives and error bounds.

\begin{definition}\label{def:theta}  \cite{LiuNgYang2009}
	The merit function $\theta : X \to \mathbb{R}\cup \left\{ +\infty \right\}$ is defined by
	 	\begin{eqnarray*}  	\theta(x) &=& -\inf_{a \in A} \Delta_C(F(x) - F(a)) = \sup_{a \in A} (-\Delta_C(F(x) - F(a))) \\
		&=&   \sup_{a \in A} \inf_{y^* \in S(C^+)} \langle y^*, F(x) - F(a) \rangle, \quad \forall x \in X.
	\end{eqnarray*}
	
	
	
\end{definition}

\begin{lemma} \label{YL:thetaprop} \cite[Theorem 3]{LiuNgYang2009}
	The function $\theta$ satisfies:
	\begin{enumerate}[leftmargin=*,label=(\roman*)]
		\item For any $x \in A$, $\theta(x) \geq 0$.
		\item $\theta$ is $C_X$-monotone, i.e.,  for any $x_1, x_2 \in X$ with $x_1 \le_{C_X} x_2$, we have 	$	  \theta(x_1)   \leq	\theta(x_2) .		$
		\item $\theta(x) = 0$ if and only if $x \in E_w$.
		\item There exists a constant $r > 0$   such that
		\[
		r \cdot d_{\widehat{A}}(x) \leq \theta(x) \leq \|F\| \cdot d_{\widehat{A}}(x), \quad \forall x \in A.
		\]
	\end{enumerate}
\end{lemma}

\begin{proposition}\label{AQWE}
	If there exists    $y_0^*\in C^+\setminus\{0\}$ such that $\inf_{a\in A}\langle y_0^*,F(a)\rangle > -\infty$, then $\theta(x)<+\infty$ for all $x\in X$.
\end{proposition}
\begin{proof}
	Suppose  that  there exists $y_0^*\in C^+\setminus\{0\}$ such that
	\[
	m:=\inf_{a\in A}\langle y_0^*,F(a)\rangle > -\infty.
	\]
	Set $\tilde y^*=y_0^*/\|y_0^*\|$.  Then $\tilde y^*\in S(C^+)$ and $\inf_{a\in A}\langle\tilde y^*,F(a)\rangle = m/\|y_0^*\| > -\infty$.   For any $x\in X$, we have
	\[
	\begin{aligned}
		\theta(x) &= \sup_{a \in A} \inf_{y \in S(C^+)} \langle y^*, F(x) - F(a) \rangle  
		\leq \sup_{a \in A} \langle \tilde y^*, F(x) - F(a) \rangle \\
		&= \langle \tilde y^*, F(x) \rangle - \inf_{a \in A} \langle \tilde y^*, F(a) \rangle  
		= \langle \tilde y^*, F(x) \rangle - m/\|y_0^*\| < + \infty.
	\end{aligned}
	\]
\end{proof}

\begin{definition}  	\cite{Luc, DF2024} 
	A nonempty set $D \subseteq Y$ is called:
	\begin{enumerate}
		\item[(i)] $C$-bounded, if for every neighborhood $U$ of $0 \in Y$, there exists $t > 0$ such that $D \subseteq tU + C$;
		\item[(ii)] $C$-sequentially compact, if for any sequence $\{y_n\}  \subseteq D$, there exist a sequence $\{c_n\}  \subseteq C$, $y_0 \in D$, and a subsequence $\{y_{n_k} - c_{n_k}\} $ of $\{y_n - c_n\} $ such that $y_{n_k} - c_{n_k} \to y_0$.
	\end{enumerate}
\end{definition}

\begin{corollary} \label{ZRUQ}
	If $F(A)$ is $C$-bounded, then $\theta(x)<+\infty$ for all $x\in X$.
\end{corollary}

\begin{proof}
	Take any $y^*\in S(C^+)$. According to the definition of $C$-boundedness, for    the open unit ball $U = \{y\in Y : \|y\| < 1\}$,  there exists $t_0 > 0$ such that
$	F(A) \subseteq t_0 U + C .$
	This means that for each $a\in A$, there exist $u_a \in U$ and $c_a \in C$ satisfying
$	F(a) = t_0 u_a + c_a .$
	Applying the functional $y^*$ to both sides yields
\begin{equation}\label{DzxcB1}
	\langle y^*, F(a)\rangle = t_0\langle y^*, u_a\rangle + \langle y^*, c_a\rangle .
 \end{equation}
	Since $y^*\in C^+$, we have $\langle y^*, c_a\rangle \ge 0$. Moreover, from $\|u_a\| < 1$ and $\|y^*\| = 1$, we get $|\langle y^*, u_a\rangle| \le \|y^*\|\cdot\|u_a\| < 1$, hence $\langle y^*, u_a\rangle > -1$. This together with  (\ref{DzxcB1})   implies that 
	\[
	\langle y^*, F(a)\rangle > -t_0 \quad \forall a\in A .
	\]
	Taking the infimum over $a$ gives
	\[
	\inf_{a\in A} \langle y^*, F(a)\rangle \ge -t_0 > -\infty .
	\]
	By Proposition \ref{AQWE}, $\theta(x)<+\infty$ for all $x\in X$.
\end{proof}

\begin{proposition}\label{XLJKD}
	If $F(A)$ is $C$-sequentially compact, then for any $x\in X$, there exists $a_0 \in A$ such that
	$$\theta(x) =  - {\Delta _C}\bigl( F(x) - F(a_0)\bigr).$$
\end{proposition}

\begin{proof}
	Fix an arbitrary $x\in X$, and set $f(y) = \Delta_C(F(x) - y)$.  It is easy to see  that $-\theta(x) = \inf_{y \in F(A)} f(y)$.
	By the $1$-Lipschitz continuity of $\Delta_C$, $f: Y \to \mathbb{R}$ is continuous.
	
	Let $y_1, y_2 \in Y$ satisfy $y_2 - y_1 \in C$, then $F(x) - y_1 - (F(x) - y_2) \in C$. By Proposition \ref{XSQ}, we have
	$$f(y_1) = {\Delta _C}(F(x) - y_1) \le {\Delta _C}(F(x) - y_2) = f(y_2).$$
	According to Proposition 3.5 of \cite{Han2025}, there exists $y_0 \in F(A)$ such that
	$${\Delta _C}(F(x) - y_0) = f(y_0) = \inf_{y \in F(A)} f(y) = -\theta(x).$$
	Since $y_0 \in F(A)$, there exists $a_0 \in A$ with $y_0 = F(a_0)$. Hence,
	$$- {\Delta _C}\bigl( F(x) - F(a_0)\bigr) = - {\Delta _C}\bigl( F(x) - y_0\bigr) = \theta(x).$$
\end{proof}

\begin{proposition} \label{QQB} 
	For a nonempty set $Q \subseteq X$, we have:
	\begin{enumerate}
		\item[(i)] if $Q$ is $C_X$-bounded, then $F(Q)$ is $C$-bounded.
		\item[(ii)] if $F$ is surjective  and   $F(Q)$ is $C$-bounded, then  $Q$ is $C_X$-bounded.
	\end{enumerate}
\end{proposition}

\begin{proof}
	(i). Take any neighborhood $V$ of $0$ in $Y$. By continuity of $F$, $U := F^{-1}(V)$ is a neighborhood of $0$ in $X$. By the $C_X$-boundedness of $Q$, there exists $t > 0$ such that $Q \subseteq tU + C_X$. Since $F$ is linear,  we have 
	\[
	F(Q) \subseteq F(tU + C_X) = tF(U) + F(C_X).
	\]
	Observe that $F(U) = F(F^{-1}(V)) \subseteq V$ and $F(C_X) \subseteq C$, hence $F(Q) \subseteq tV + C$.
	By the arbitrariness of $V$, $F(Q)$ is $C$-bounded.
	
	(ii). Since $X, Y$ are Banach spaces and $F$ is a continuous linear surjection, by the Banach--Schauder open mapping theorem, $F$ is an open map. Take any open neighborhood $U$ of $0$ in $X$. Set $V := F(U)$. By the open mapping property, $V$ is an open set in $Y$ containing $0 \in Y$. By the $C$-boundedness of $F(Q)$, there exists $t > 0$ such that $F(Q) \subseteq tV + C$. For any $x \in Q$, since $F(x) \in F(Q) \subseteq tV + C$, there exist $v \in V$ and $c \in C$ such that $F(x) = t v + c$. From $V = F(U)$, there exists $u \in U$ satisfying $v = F(u)$. Thus,
	\[
	F(x) = t F(u) + c \implies F(x - t u) = c \in C.
	\]
	By $C_X = F^{-1}(C)$, we have $x - t u \in C_X$. Hence $x = t u + (x - t u) \in tU + C_X$.
	By the arbitrariness of $x$, $Q \subseteq tU + C_X$. Therefore, $Q$ is $C_X$-bounded.
\end{proof}


\begin{remark}  By   Corollary  \ref{ZRUQ}  and Proposition  \ref{QQB}  (i), we get that  if $A$ is $C_X$-bounded, then $\theta(x)<+\infty$ for all $x\in X$.
   \end{remark}


\begin{remark}   \label{VKYJQ}  Let $y^* \in  C^+$ and  $A$ be $C_X$-bounded.  It follows from  Proposition  \ref{QQB}  (i)  that $F(A)$ is $C$-bounded. Then it is easy to obtain that 
	$$\mathop {\inf }\limits_{a \in A} \left\langle {{F^*}{y^*},a} \right\rangle  = \mathop {\inf }\limits_{a \in A} \left\langle {{y^*},F\left( a \right)} \right\rangle  >  - \infty $$
	
  \end{remark}


\begin{lemma}\label{XDRT}
	If $F$ is surjective, then there exists a constant $\lambda>0$ such that for every $x\in X$,
	\[
	\operatorname{dist}(x, C_X) \le \lambda \operatorname{dist}(F(x), C).
	\]
\end{lemma}

\begin{proof}
	Let $\delta = \operatorname{dist}(F(x), C)$. 
	If \(\delta=0\),  it follows from   the closedness of  \(C\)   that 
	 \(F(x)\in C\), and hence \(x\in C_X\). Therefore, the result is trivial. In the following, we consider the case \(\delta>0\).
	
	
	
	
	For any $\varepsilon > 0$,   there exists $c_\varepsilon \in C$ such that
	\begin{equation}\label{eq:1}
		\|F(x) - c_\varepsilon\| < \delta + \varepsilon .
	\end{equation}
	Since $F$ is a continuous linear surjection, the Banach--Schauder open mapping theorem asserts that $F$ is open. Hence there exists $r>0$ such that
	\begin{equation}\label{eq:2}
		r B_Y \subseteq F(B_X).
	\end{equation}
 
	
	Let $y_\varepsilon := F(x) - c_\varepsilon \in Y$.	We conclude from  \(\delta>0\)  that    $y_\varepsilon \neq 0$,  and define
$	y_r := r \frac{y_\varepsilon}{\|y_\varepsilon\|}.$
	Then $y_r \in r B_Y$. By \eqref{eq:2}, there exists $v_r \in B_X$ such that $F(v_r) = y_r$. Let
$	u_r := \frac{\|y_\varepsilon\|}{r} v_r.$
	By linearity of $F$,
	\begin{equation}\label{eq:++}
	F(u_r) = \frac{\|y_\varepsilon\|}{r} F(v_r) = \frac{\|y_\varepsilon\|}{r} \cdot r \frac{y_\varepsilon}{\|y_\varepsilon\|} = y_\varepsilon.
	\end{equation}
	Moreover,
	\begin{equation}\label{eq:3}
		\|u_r\| = \frac{\|y_\varepsilon\|}{r} \|v_r\| \leq \frac{\|y_\varepsilon\|}{r} .
	\end{equation}
It follows from  (\ref{eq:++})  that  $	F(x - u_r) = F(x) - y_\varepsilon = c_\varepsilon \in C,$
and	so $x - u_r \in F^{-1}(C) = C_X$. Combining \eqref{eq:3} and $y_\varepsilon = F(x) - c_\varepsilon$, we obtain
	\begin{equation}\label{eq:4}
		\operatorname{dist}(x, C_X) \le \|x - (x - u_r)\| = \|u_r\| \leq \frac{\|y_\varepsilon\|}{r} = \frac{1}{r} \|F(x) - c_\varepsilon\|.
	\end{equation}
	Together with \eqref{eq:1}, we have
$	\operatorname{dist}(x, C_X) \le \frac{1}{r} (\delta + \varepsilon).$
	Since $\varepsilon > 0$ is arbitrary,  we get
	\[
	\operatorname{dist}(x, C_X) \le \frac{1}{r} \delta = \frac{1}{r} \operatorname{dist}(F(x), C).
	\]
\end{proof}

\begin{proposition}
	For a nonempty set $Q \subseteq X$, we have:
	\begin{enumerate}
		\item[(i)] if $Q$ is $C_X$-sequentially compact, then $F(Q)$ is $C$-sequentially compact.
		\item[(ii)] if $F$ is surjective     and    $F(Q)$ is $C$-sequentially compact, then   $Q$ is $C_X$-sequentially compact.
	\end{enumerate}
\end{proposition}

\begin{proof}
	(i). This follows from Lemma 2.3 of   \cite{Han2026}.
	
	(ii). Let $\{x_n\} \subseteq Q$ be an arbitrary sequence. Set $y_n:=F(x_n)\in F(Q)$. By $C$-sequential compactness of $F(Q)$, there exist a sequence $\{c_n\}  \subseteq C$, $y_0\in F(Q)$, and a subsequence $\{y_{n_k}\} $ of  $\{y_n\}$ such that
$	y_{n_k}-c_{n_k}\to y_0  $.
	Since $y_0\in F(Q)$, there exists $x_0\in Q$ with $F(x_0)=y_0$. Consider $u_k:=x_{n_k}-x_0\in X$. Then
	\[
	F(u_k)=F(x_{n_k})-F(x_0)=y_{n_k}-y_0.
	\]
	From $y_{n_k}-c_{n_k}\to y_0$, we have $y_{n_k}-y_0-c_{n_k}\to0$, i.e.,
$	F(u_k)-c_{n_k}\to 0 .$
	Because $c_{n_k}\in C$, we obtain $\dist(F(u_k),C)\le \|F(u_k)-c_{n_k}\|\to0$.
	Applying Lemma \ref{XDRT} to each $u_k$, we  have 
	\[
	\dist(u_k,\,C_X)\le \lambda \,\dist(F(u_k),C) \to 0 \qquad(k\to\infty).
	\]
Then for each $k$, there exists $w_k\in C_X$ such that $\|u_k-w_k\|<\dist(u_k,C_X)+\frac1k$. Hence $\|u_k-w_k\|\to0$.
	This together with  $u_k=x_{n_k}-x_0$  implies that 
	$	x_{n_k}-w_k = x_0 + u_k - w_k \to x_0 .$
	Set $d_{n_k}:=w_k\in C_X$ (for other indices one may define arbitrarily, e.g., $d_n=0$ if $n\notin\{n_k\}$). Then the sequence $\{d_n\}\subseteq C_X$ satisfies $x_{n_k}-d_{n_k}\to x_0$. Therefore,  $Q$ is $C_X$-sequentially compact.
\end{proof}

For the convenience of discussion, from now on we always assume that $\theta(x)<+\infty$ for all $x\in X$. The following two properties of $\theta$ are important because they play a key role in the derivation of several results in this paper.

\begin{theorem}\label{YL:thetaprop2}
	The  merit function $\theta$ satisfies:
	\begin{enumerate}
		\item[(i)] $\theta$ is concave on $X$;
		\item[(ii)] $\theta$ is Lipschitz continuous on $X$ with constant $\|F\|$.
	\end{enumerate}
\end{theorem}

\begin{proof}
	(i). Since $C$ is convex, $\Delta_C$ is convex, and therefore $-\Delta_C$ is concave on $Y$. Define $g: X\times A \to \mathbb{R}$ by
	\[
	g(x,a) = -\Delta_C\big(F(x)-F(a)\big),  \quad  \forall  (x,a)  \in  X\times A.
	\]
	For any $(x_1,a_1),(x_2,a_2) \in X\times A$ and $\lambda\in[0,1]$, we have
	\begin{eqnarray} \label{XSRT1}
		g\bigl( \lambda x_1 + (1-\lambda)x_2, \lambda a_1 + (1-\lambda)a_2 \bigr)
		&=& -\Delta_C\big( \lambda(F(x_1)-F(a_1)) + (1-\lambda)(F(x_2)-F(a_2)) \big) \nonumber\\ 
		&\ge& \lambda\big[-\Delta_C(F(x_1)-F(a_1))\big] + (1-\lambda)\big[-\Delta_C(F(x_2)-F(a_2))\big] \nonumber\\ 
		&=& \lambda g(x_1,a_1) + (1-\lambda)g(x_2,a_2).  
\end{eqnarray}
It is clear that
	\begin{equation} \label{XSRT2}
		\theta(x) = -\inf_{a\in A} \Delta_C\big(F(x)-F(a)\big)
		= \sup_{a\in A} \Big[-\Delta_C\big(F(x)-F(a)\big)\Big]
		= \sup_{a\in A} g(x,a), \quad \forall x \in X.  
	\end{equation}
Let $x_1, x_2 \in X$ and $\lambda\in[0,1]$.
	For any $\varepsilon>0$, by (\ref{XSRT2}), there exist $a_1,a_2\in A$ such that
\begin{equation} \label{XSRT+2}
	g(x_1,a_1) > \theta(x_1) - \varepsilon, \qquad g(x_2,a_2) > \theta(x_2) - \varepsilon.
\end{equation}
	Set $a_\lambda = \lambda a_1 + (1-\lambda)a_2$. Since $A$ is convex, $a_\lambda \in A$. Then  it follows from  (\ref{XSRT2})  that
\begin{equation} \label{XSRT3}
	\theta(\lambda x_1 + (1-\lambda)x_2) \ge g\big(\lambda x_1 + (1-\lambda)x_2,\; a_\lambda\big).
	\end{equation}
    We conclude from  (\ref{XSRT1})  and  (\ref{XSRT+2})  that
	\[
	\begin{aligned}
		g\big(\lambda x_1 + (1-\lambda)x_2,\; a_\lambda\big)
		&\ge \lambda g(x_1,a_1) + (1-\lambda)g(x_2,a_2) \\
		&> \lambda(\theta(x_1)-\varepsilon) + (1-\lambda)(\theta(x_2)-\varepsilon) \\
		&= \lambda\theta(x_1) + (1-\lambda)\theta(x_2) - \varepsilon.
	\end{aligned}
	\]
	  This together with   (\ref{XSRT3})  implies that 
	\[
	\theta(\lambda x_1 + (1-\lambda)x_2) \ge \lambda\theta(x_1) + (1-\lambda)\theta(x_2) - \varepsilon.
	\]
	By the arbitrariness of $\varepsilon>0$, we obtain
$
	\theta(\lambda x_1 + (1-\lambda)x_2) \ge \lambda\theta(x_1) + (1-\lambda)\theta(x_2).
$
	
	(ii). Fix arbitrary $x_1,x_2\in X$. For any given $\epsilon>0$, there exists $a_\epsilon\in A$ such that
	\begin{equation}   \label{XSRQT1}
		\theta(x_1)-\epsilon < \inf_{y^*\in S(C^+)}\langle y^*,\,F(x_1)-F(a_\epsilon)\rangle .  
	\end{equation}
	For this $a_\epsilon\in A$, we get 
$$
		\theta(x_2)=\sup_{a\in A}\inf_{y^*\in S(C^+)}\langle y^*,\,F(x_2)-F(a)\rangle
		\ge \inf_{y^*\in S(C^+)}\langle y^*,\,F(x_2)-F(a_\epsilon)\rangle , 
$$
 and so
	\begin{equation}\label{XSRQT2}
		-\theta(x_2)\le -\inf_{y^*\in S(C^+)}\langle y^*,\,F(x_2)-F(a_\epsilon)\rangle .  
	\end{equation}
	Adding  (\ref{XSRQT1})  and  (\ref{XSRQT2}), we obtain
	\begin{align*}
		\theta(x_1)-\theta(x_2)-\epsilon 
		&< \inf_{y^*\in S(C^+)}\langle y^*,\,F(x_1)-F(a_\epsilon)\rangle
		-\inf_{y^*\in S(C^+)}\langle y^*,\,F(x_2)-F(a_\epsilon)\rangle \\
		&\le \sup_{y^*\in S(C^+)}\Big[
		\langle y^*,\,F(x_1)-F(a_\epsilon)\rangle
		-\langle y^*,\,F(x_2)-F(a_\epsilon)\rangle
		\Big] \\
		&= \sup_{y^*\in S(C^+)} \langle y^*,\,F(x_1-x_2)\rangle  
		\le \sup_{y^*\in S(C^+)} \|y^*\|\,\|F(x_1-x_2)\| \\
		&= \|F(x_1-x_2)\| \le \|F\|\,\|x_1-x_2\|.
	\end{align*}
	Since $\epsilon>0$ is arbitrary, we have
$	\theta(x_1)-\theta(x_2) \le \|F\|\,\|x_1-x_2\|. $
	By symmetry,  we get
$	\theta(x_2)-\theta(x_1) \le \|F\|\,\|x_1-x_2\|.$
	Hence,
	$	|\theta(x_1)-\theta(x_2)| \le \|F\|\,\|x_1-x_2\|.$
\end{proof}

\begin{theorem}  \label{TxEQ} 
	Let $x\in X$. The following assertions hold:
	\begin{enumerate}
		\item[(i)] $x\in A + \operatorname{int}C_X$ if and only if $\theta(x)>0$;
		\item[(ii)] If $x\in A + C_X$, then $\theta(x)\ge 0$;
		\item[(iii)] If $\theta(x)\ge 0$ and $F(A)$ is $C$-sequentially compact, then $x\in A + C_X$.
	\end{enumerate}
\end{theorem}

\begin{proof}
	For any $x\in X$, by definition of $\theta$,
	\begin{equation}\label{eq:theta-dist}
		\theta(x)=\sup_{a\in A}\bigl[-\Delta_C(F(x)-F(a))\bigr].
	\end{equation}
	
	\textbf{(i)} Suppose $x\in A+\operatorname{int}C_X$. Then there exist $a_0\in A$ and $d\in\operatorname{int}C_X$ such that $x=a_0+d$. By   $\operatorname{int}C_X=F^{-1}(\operatorname{int}C)$, we get $F(d)\in\operatorname{int}C$. It follows from   Proposition \ref{XSQ} that
	$${\Delta _C}\bigl( F(x) - F(a_0) \bigr) = {\Delta _C}\bigl( F(x - a_0) \bigr) = {\Delta _C}\bigl( F(d) \bigr) < 0.$$
	By \eqref{eq:theta-dist}, we have  $\theta(x) \ge -{\Delta _C}\bigl( F(x) - F(a_0) \bigr) > 0$.
	
	Conversely, suppose $\theta(x)>0$. By \eqref{eq:theta-dist}, there exists $a' \in A$ such that
	$-\Delta_C(F(x) - F(a')) > 0$.
	By Proposition \ref{XSQ},
	$F(x) - F(a') \in \operatorname{int} C$.
	Hence
	$x - a' \in F^{-1}(\operatorname{int} C) = \operatorname{int} C_X$.
	Therefore
	$$x \in a' + \operatorname{int} C_X \subseteq A + \operatorname{int} C_X.$$
	
	\textbf{(ii)} Suppose $x\in A + C_X$. Then there exist $a_0\in A$ and $d\in C_X$ such that $x=a_0+d$. Hence
	\[
	F(x)-F(a_0)=F(d)\in F(C_X)\subseteq C.
	\]
	Thus ${\Delta _C}\bigl( F(x) - F(a_0) \bigr) \le0$.     It follows from  (\ref{eq:theta-dist})  that
	$\theta(x) \ge -{\Delta _C}\bigl( F(x) - F(a_0) \bigr) \ge 0$.
	
	\textbf{(iii)} Since $F(A)$ is $C$-sequentially compact and $\theta(x)\ge 0$, by Proposition \ref{XLJKD}, there exists $a_0 \in A$ such that
	$$\theta(x) = -{\Delta _C}\bigl( F(x) - F(a_0) \bigr) \ge 0.$$
	By Proposition \ref{XSQ} and the closedness of $C$, we have $F(x)-F(a_0) \in C$, hence $x - a_0 \in C_X$, and therefore
	$x \in a_0 + C_X \subseteq A + C_X.$
\end{proof}

 

\begin{lemma}[Sion's Minimax Theorem]\label{SJDJX}
	Let $H$ be a nonempty compact convex set in a Hausdorff topological vector space, and $Q$ a nonempty convex set in a topological vector space. Suppose the function $f : H \times Q \to \mathbb{R}$ satisfies:
	\begin{enumerate}
		\item For each $y \in Q$, $x \mapsto f(x,y)$ is lower semicontinuous and convex on $H$;
		\item For each $x \in H$, $y \mapsto f(x,y)$ is concave on $Q$.
	\end{enumerate}
	Then
	\[
	\min_{x \in H} \sup_{y \in Q} f(x,y) = \sup_{y \in Q} \min_{x \in H} f(x,y).
	\]
\end{lemma}

\begin{theorem}[Dual representation of the merit function]\label{YDOBS}
	The merit function $\theta$ can be represented as:
	\[
	\theta(x)=\min_{v\in K}\Bigl(\langle v,F(x)\rangle -\inf_{a\in A}\langle v,F(a)\rangle\Bigr)
	=\min_{v\in K}\Bigl(\langle F^*v, x\rangle+\sigma_A(-F^*v)\Bigr), \quad \forall x\in X.
	\]
	where $K=\overline{\operatorname{co}}^{w^*}\bigl(S(C^+)\bigr)$ and $\sigma_A(u)=\sup_{a\in A}\langle u,a\rangle$ is the support function of $A$.
\end{theorem}

\begin{proof}	For any fixed $z\in Y$, the linear functional $y^*\mapsto\langle y^*,z\rangle$ is weak$^*$-continuous on $Y^*$, hence attains its minimum on the weak$^*$-compact set $K$, and this minimum equals the infimum over $S(C^+)$, 
\begin{equation}\label{DmlB1}
	\inf_{y^*\in S(C^+)}\langle y^*,z\rangle = \min_{v\in K}\langle v,z\rangle .  
 \end{equation}
	By definition of $\theta$ and  (\ref{DmlB1}),
\begin{equation}\label{DmlB2}
	\theta(x) = \sup_{a\in A} \inf_{y^*\in S(C^+)} \langle y^*, F(x)-F(a)\rangle
	= \sup_{a\in A}\, \min_{v\in K} \langle v, F(x)-F(a)\rangle .  
 \end{equation}
	Define the bifunction $\varphi : A\times K  \to  \mathbb{R} $ by
	\[
	\varphi(a,v) = \langle v, F(x)-F(a)\rangle , \quad   \forall   (a,v)\in A\times K .
	\]
 	By Lemma \ref{SJDJX} (Sion's minimax theorem), we have
\begin{equation}\label{DmlB3}
	\sup_{a\in A}\, \min_{v\in K} \varphi(a,v) = \min_{v\in K}\, \sup_{a\in A} \varphi(a,v).  
\end{equation}
	Then,
	\begin{eqnarray} \label{DmlB4}
		\sup_{a\in A}\varphi(a,v) &=& \sup_{a\in A}\langle v, F(x)-F(a)\rangle \nonumber\\ 
		&=& \langle v, F(x)\rangle + \sup_{a\in A}\langle -v, F(a)\rangle  		 = \langle v, F(x)\rangle - \inf_{a\in A}\langle v, F(a)\rangle .
\end{eqnarray}
	Using the adjoint operator $F^*:Y^*\to X^*$, we have
	$$\langle v, F(x)\rangle = \langle F^*v, x\rangle, \quad \forall x\in X,$$
	and
	\[
	-\inf_{a\in A}\langle v, F(a)\rangle = \sup_{a\in A}\langle -F^*v, a\rangle = \sigma_A(-F^*v).
	\]
	Combining (\ref{DmlB2}), (\ref{DmlB3})  and  (\ref{DmlB4}), we obtain
	\[
	\theta(x) = \min_{v\in K} \Bigl( \langle v, F(x)\rangle - \inf_{a\in A}\langle v, F(a)\rangle \Bigr)
	= \min_{v\in K} \Bigl( \langle F^*v, x\rangle + \sigma_A(-F^*v) \Bigr).
	\]
\end{proof}

 

From  Theorem 2.4.18 (p.~97) of \cite{Zalinescu2002}, we obtain the following lemma.

 



\begin{lemma}\label{ITDL}
	Let   $X$ be a locally convex Hausdorff topological vector space,  $Q$ be a nonempty compact Hausdorff space,  $g: Q \to (X^*, w^*)$ be continuous with respect to the weak$^*$-topology on $X^*$,  and $s: Q \to \mathbb{R}$ be upper semicontinuous.
 	For each $v \in Q$, define the affine function
	\[
	f_v(x) = \langle g(v), x \rangle + s(v), \qquad x \in X.
	\]
	Set
	\[
	\phi(x) = \max_{v \in Q} f_v(x), \qquad x \in X.
	\]
	Then $\phi: X \to \mathbb{R}$ is a convex function.  If $\phi$ is continuous at some point $\bar x \in X$, then its subdifferential $\partial \phi(\bar x)$ is nonempty and satisfies
	\[
	\partial \phi(\bar x) = \overline{\operatorname{co}}^{w^*} \{ g(v) : v \in Q(\bar x) \},
	\]
	where $Q(\bar x) = \{ v \in Q : f_v(\bar x) = \phi(\bar x) \}$ is the active index set.
\end{lemma}




Let $\xi = -\theta$. Clearly $\xi$ is a convex   function on $X$.  Denote by $\partial \xi({x})$ the subdifferential of the convex function $\xi$ at $x \in X $, i.e.,
\[
\partial \xi(\bar{x}) := \big\{ x^* \in X^* : \xi(y) \ge \xi(\bar{x}) + \langle x^*, y - x \rangle,\ \forall y \in X \big\}.
\]

For the concave function \(\theta\), the superdifferential at a point \(x\in X\) is defined by
\[
\partial^+ \theta(x):=\{x^*\in X^*:\theta(y)\le\theta(x)+\langle x^*,y-x\rangle, \; \forall y\in X\}.
\]

\begin{theorem}  \label{BSXEQ}  Let $A  \subseteq X$ be $C_X$-bounded.
 For any $\bar{x} \in E_w$, the following representation holds:
	\[
	-\partial \xi(\bar{x}) = \overline{\operatorname{co}}^{w^*} \bigl\{ F^*v : v \in K,\; F^*v \in -N_A(\bar{x}) \bigr\} = \bigl( -N_A(\bar{x}) \bigr) \cap \overline{\operatorname{co}}^{w^*} \bigl( F^*(S(C^+)) \bigr) = \bigl( -N_A(\bar{x}) \bigr) \cap F^*(K).
	\]
\end{theorem}

\begin{proof}
	By Theorem \ref{YDOBS}, we have
$$
	\theta(x) = \min_{v \in K} \bigl( \langle F^*v, x \rangle + \sigma_A(-F^*v) \bigr),  
$$
 and so 
\begin{equation}\label{XCB1}
	\xi(x) = -\theta(x) = \max_{v \in K} \bigl( -\langle F^*v, x \rangle - \sigma_A(-F^*v) \bigr).  
 \end{equation}
	For each $v \in K$, define
	\[
	h_v(x) := \langle -F^*v, x \rangle - \sigma_A(-F^*v), \quad x \in X.
	\]
	Then,  it follows from  (\ref{XCB1})  that
	\[
	\xi(x) = \max_{v \in K} h_v(x).
	\]
	
	For $\bar{x} \in E_w$, we have $\theta(\bar{x}) = 0$, hence $\xi(\bar{x}) = 0$. Define the active index set
	\[
	K(\bar{x}) := \{ v \in K : h_v(\bar{x}) = \xi(\bar{x}) = 0 \}.
	\]
		For any $v \in K$, set $g(v) = -F^*v$ and $s(v) = - \sigma_A(-F^*v)$.   Since $A$ is $C_X$-bounded, it follows from  $K \subseteq C^+$ and   Remark   \ref{VKYJQ}    that $s(v) \in \mathbb{R}$ for any $v \in K$.
		  It is readily seen that   $g$ is continuous   and   $s$ is upper semicontinuous with respect to the weak$^*$-topology.
		By Lemma \ref{ITDL}, we have
$	\partial \xi(\bar{x}) = \overline{\operatorname{co}}^{w^*} \bigl\{ -F^*v : v \in K(\bar{x}) \bigr\}, $
and so 
\begin{equation}\label{XCB2}
	- \partial \xi(\bar{x}) = \overline{\operatorname{co}}^{w^*} \bigl\{ F^*v : v \in K(\bar{x}) \bigr\}.  
 \end{equation}
	
	By definition of $K(\bar{x})$, $v \in K(\bar{x})$ iff $v \in K$ and
\begin{equation}\label{XCB3}
	0 = \langle F^*v, \bar{x} \rangle + \sigma_A(-F^*v).  
 \end{equation}
	The support function gives
	\[
	\sigma_A(-F^*v) = \sup_{a \in A} \langle -F^*v, a \rangle = - \inf_{a \in A} \langle F^*v, a \rangle.
	\]
	Substituting into  (\ref{XCB3}) yields
\begin{equation}\label{XCB4}
	0 = \langle F^*v, \bar{x} \rangle - \inf_{a \in A} \langle F^*v, a \rangle
	\quad \Longleftrightarrow \quad
	\langle F^*v, \bar{x} \rangle = \inf_{a \in A} \langle F^*v, a \rangle.  
 \end{equation}
		Equation (\ref{XCB4}) means that $\bar{x}$ is a global minimizer of the continuous linear functional $a \mapsto \langle F^*v, a \rangle$ on the closed convex set $A$. This is equivalent to
$	-F^*v \in N_A(\bar{x}).$
	Combining with (\ref{XCB2}), we have
\begin{equation}\label{XCB++}
	-\partial \xi(\bar{x}) = \overline{\operatorname{co}}^{w^*} \bigl\{ F^*v : v \in K,\; F^*v \in -N_A(\bar{x}) \bigr\}.
 \end{equation}
	
	Next, we prove that
\begin{equation}\label{XCB5}
	F^*(K) = \overline{\operatorname{co}}^{w^*}\bigl( F^*(S(C^{+})) \bigr).  
 \end{equation}
	First we show $F^*(K) \subseteq \overline{\operatorname{co}}^{w^*}\bigl( F^*(S(C^{+})) \bigr)$.
In fact,	since $F^*$ is linear, for any finite convex combination, we have
	\[
	F^*\left( \sum_{i=1}^n \lambda_i y_i^* \right) = \sum_{i=1}^n \lambda_i F^*(y_i^*), \quad \lambda_i \ge 0,\ \sum_{i=1}^n \lambda_i = 1, \;  y_i^* \in S(C^+).
	\]
	This means
\begin{equation}\label{XCB6}
	F^*(\operatorname{co}(S(C^+))) = \operatorname{co}( F^*(S(C^+)) ).  
 \end{equation}
	Take any $z \in K = \overline{\operatorname{co}}^{w^*}(S(C^+))$. By definition of the weak$^*$-closure, there exists a net $\{z_\alpha\}_{\alpha \in I} \subseteq \operatorname{co}(S(C^+))$ such that $z_\alpha \xrightarrow{w^*} z$.
	For each $\alpha$, by (\ref{XCB6}) we have $F^*(z_\alpha) \in \operatorname{co}(F^*(S(C^+)))$. Since $F^*$ is weak$^*$-to-weak$^*$ continuous, we get 
	\[
	F^*(z_\alpha) \xrightarrow{w^*} F^*(z) \quad \text{in } X^*.
	\]
 This means that
$	F^*(z) \in \overline{\operatorname{co}}^{w^*}\bigl( F^*(S(C^+)) \bigr).$
	Since $z \in K$ is arbitrary, we obtain 
	\begin{equation}\label{XCB7}
	F^*(K) \subseteq \overline{\operatorname{co}}^{w^*}\bigl( F^*(S(C^+)) \bigr).
	 \end{equation}
	 
	Since $K = \overline{\operatorname{co}}^{w^*}(S(C^+))$ is   a   weak$^*$-compact convex set, and $F^*$ is   linear and weak$^*$-to-weak$^*$ continuous, the set $F^*(K)$ is weak$^*$-compact and convex. This together with $F^*(S(C^+)) \subseteq F^*(K)$  implies that
$	\overline{\operatorname{co}}^{w^*} \bigl( F^*(S(C^+)) \bigr) \subseteq F^*(K).$   Combining this with (\ref{XCB7}), we get that (\ref{XCB5})  holds.
	
 
	Since $K = \overline{\operatorname{co}}^{w^*}(S(C^{+}))$, and $K(\bar{x}) \subseteq K$, using (\ref{XCB5}),
	\[
	\{ F^*v : v \in K(\bar{x}) \} \subseteq F^*(K) = \overline{\operatorname{co}}^{w^*}\bigl( F^*(S(C^{+})) \bigr).
	\]
	From (\ref{XCB2}) we obtain
		\begin{equation}\label{XCB8}
	- \partial \xi(\bar{x}) \subseteq \overline{\operatorname{co}}^{w^*} \bigl( F^*(S(C^{+})) \bigr).  
	 \end{equation}
	
Noting that   $-N_A(\bar{x})$ is convex and weak$^*$-closed, 	it is easy to see that $ \overline{\operatorname{co}}^{w^*} \bigl\{ F^*v : v \in K,\; F^*v \in -N_A(\bar{x}) \bigr\}    \subseteq -N_A(\bar{x})$. 	This together with  (\ref{XCB++})  yields that   $-\partial \xi(\bar{x}) \subseteq -N_A(\bar{x})$.
		Combining with (\ref{XCB8}), we have
	\[
	- \partial \xi(\bar{x}) \subseteq \bigl( -N_A(\bar{x}) \bigr) \cap \overline{\operatorname{co}}^{w^*} \bigl( F^*(S(C^{+})) \bigr).
	\]
	
 
	Take any $x^* \in \bigl( -N_A(\bar{x}) \bigr) \cap \overline{\operatorname{co}}^{w^*} \bigl( F^*(S(C^+)) \bigr)$.
	By (\ref{XCB5}), we get  $x^* \in F^*(K)$, so there exists $v \in K$ with $x^* = F^*v$.
	Since $F^*v  =    x^* \in -N_A(\bar{x})$,     we have $x^* \in \bigl\{ F^*v : v \in K,\; F^*v \in -N_A(\bar{x}) \bigr\}$.   This together with  (\ref{XCB++})  implies that  $x^* \in   -\partial \xi(\bar{x}).$
	 	This proves
	\[
	\bigl( -N_A(\bar{x}) \bigr) \cap \overline{\operatorname{co}}^{w^*} \bigl( F^*(S(C^+)) \bigr) \subseteq -\partial \xi(\bar{x}).
	\]
\end{proof}

\begin{remark}  \label{RXExcQ}  
	It follows from   $\xi = -\theta$  that  $	\partial^+ \theta(x) =  	- \partial \xi({x}) $.   Let $A  \subseteq X$ be $C_X$-bounded.    Then by Theorem \ref{BSXEQ}, we obtain that if $\bar{x} \in E_w$, then 
		\[
		\partial^+ \theta(x) = \overline{\operatorname{co}}^{w^*} \bigl\{ F^*v : v \in K,\; F^*v \in -N_A(\bar{x}) \bigr\} = \bigl( -N_A(\bar{x}) \bigr) \cap \overline{\operatorname{co}}^{w^*} \bigl( F^*(S(C^+)) \bigr) = \bigl( -N_A(\bar{x}) \bigr) \cap F^*(K).
	\]
	
	
	
  \end{remark}



\section{Error bounds   of the merit   function  }\label{sec:errorbounds}

 In this section we undertake a comprehensive study of error bounds for \(\theta\). We establish the equivalence of thirteen different characterizations, ranging from linear regularity of the pair \(\{A,\widehat A\}\), through  the  global slope condition, to perturbation stability and sublevel‐set Hausdorff continuity. These equivalences unify and extend previous results in the literature and set the stage for the directional‐derivative characterizations developed in later sections.  

\begin{lemma}[Ekeland's variational principle]\label{ELdaz}
	Let $(M, d)$ be a complete metric space and $f: M \to \mathbb{R} \cup \{+\infty\}$ a proper lower semicontinuous function bounded below. If $x_0 \in \operatorname{dom} f$ and $\varepsilon > 0,\ \lambda > 0$ satisfy
	\[
	f(x_0) \le \inf_{x \in M} f(x) + \varepsilon,
	\]
	then there exists $\bar{x} \in M$ such that:
	\begin{enumerate}
		\item $f(\bar{x}) \le f(x_0)$;
		\item $d(x_0, \bar{x}) \le \lambda$;
		\item For all $x \in M \setminus \{\bar{x}\}$,
		\[
		f(x) > f(\bar{x}) - \frac{\varepsilon}{\lambda} \, d(x, \bar{x}) .
		\]
	\end{enumerate}
\end{lemma}



 



The set-valued map $\Phi:[0,+\infty)\rightrightarrows A$ is defined by
$$\Phi(t)=\{x\in A:\theta(x)\le t\}, \quad  \forall t \in [0,+\infty).$$

For any \(x\in A\) with \(\theta(x)>0\),  the  global slope is defined by
\[
m(x):=\sup_{\substack{y\in A\\ \theta(y)<\theta(x)}}\frac{\theta(x)-\theta(y)}{\|x-y\|}\in[0,+\infty].
\]


\begin{theorem} \label{XLWCJK}  Consider the following statements:
	\begin{enumerate}
		\item[(i)] $\theta$ has an error bound on $A$: there exists $\tau>0$ such that
		\begin{equation}\label{EBAQJ} 
			d(x,E_w)\le\tau \theta(x), \quad  \forall x\in A.
		\end{equation}		 
		\item[(ii)] $\{A,\widehat{A}\}$ is linearly regular: there exists $\kappa>0$ such that
				 $$d_{A\cap\widehat{A}}(x)\le\kappa\max\{d_A(x),d_{\widehat{A}}(x)\},   \quad  \forall     x\in X. $$  
		 
		\item[(iii)] $d_{\widehat{A}}$ has an error bound on $A$: there exists $\tau>0$ such that $d(x,E_w)\le\tau d_{\widehat{A}}(x)$ for all $x\in A$.
		\item[(iv)]  There exist  \(\alpha > 0\) and \(\beta \in (0,1)\) such that for every \(x \in A\) one can find \(y \in A\) satisfying
		\begin{equation}\label{EYZXJ} 
		\|x - y\| \le \alpha\,\theta(x) \quad\text{and}\quad \theta(y) \le \beta\,\theta(x).
		\end{equation}	
		\item[(v)]  There exist \(\kappa > 0\) such that  
		\begin{equation}\label{EBkkQJ} 
		d(x+e, E_w) \le \kappa (\theta(x) + \|e\|),  \quad  \forall x \in A, \;    \forall e \in X.
			\end{equation}	
		
		\item[(vi)]  There exist $\gamma>0$ such that  
		\begin{equation}\label{EBkXJ} 
		d(x,E_w)\le \gamma\bigl(\max\{\theta(x),0\}+d_A(x)\bigr),\quad \forall x\in X .
				\end{equation}	
		
		\item[(vii)]  There exists   \(\kappa > 0\) such that  
			$$			d(x, E_w) \le \kappa \, \zeta(x), \quad \forall x \in X, 		$$
			where  $\zeta(x) := \inf_{a \in A} \bigl( \theta(a) + \|x - a\| \bigr)$.
			
			
			
			\item[(viii)] There exists   \(c>0\) such that for each \(x\in A\setminus E_w\), one can find \(y\in A\) satisfying
			\begin{equation}\label{EevL}
				y\neq x \quad\text{and}\quad
				\theta(y)\le \theta(x)-c\|x-y\| .
	 \end{equation}
			
					\item[(ix)]   There exists   \(\gamma >0\) such that
				\[
				m(x)\ge \gamma,\quad\forall x\in A\setminus E_w.
				\]
			
			 	\item[(x)]  There exists $c>0$ such that    
 	$$ d \bigl(F(x),\operatorname{WMin}(F(A),C)\bigr)     \le c \; \theta(x),   \quad  \forall x\in A.$$ 
		
		\item[(xi)]  Asymptotic   condition:
		\[
		\liminf_{\substack{x\xrightarrow{A} E_w\\ x\notin E_w}}\frac{d_{\widehat{A}}(x)}{d(x,E_w)}>0.
		\]
		where $x\xrightarrow{A} E_w$ means $x\in A$ and $d(x,E_w)\to 0$.
		
		\item[(xii)] There exist \(t_0 > 0\) and \(\beta > 0\) such that 
		\[
		d_H(\Phi(t), E_w) \le \beta t,  \quad  \forall t \in [0, t_0].
		\]
 
		
	\item[(xiii)] There exist $\kappa>0$ and an open set $U\supseteq E_w$ such that \begin{equation}\label{EBkmJ} 
			d(x,\Phi(0))\le\kappa\,d(0,\Phi^{-1}(x)),       \quad  \forall       x\in U\cap A. 	
			\end{equation}	 
	\end{enumerate}
	 	Then the following hold:
		\begin{enumerate}
		\item The statements (i)--(ix) are equivalent, (i) $\Rightarrow$ (x), (iii) $\Rightarrow$ (xi),  (i) $\Rightarrow$ (xii) and (i) $\Rightarrow$ (xiii). 
			\item  If there exist $\mu>0$ such that $\|F(z)\|\ge\mu\|z\|$ for any $z\in A-A$,
	then (x) $\Rightarrow$ (i).
			\item  If $A$ is bounded, then (xii) $\Rightarrow$ (i).
	 \item  	If $A$ is compact, then (xi) $\Rightarrow$ (iii) and (xiii) $\Rightarrow$ (i).  
	\end{enumerate}
\end{theorem}

\begin{proof}
It follows from Lemma \ref{YL:thetaprop}     that  (i) $\Leftrightarrow$ (iii).
		By \cite[Theorem 4]{LiuNgYang2009},  we get    (ii) $\Leftrightarrow$ (iii).
	
 
	\noindent\textbf{(i) $\Rightarrow$ (iv)}.
  Assume  that   there exists $\tau>0$ such that  (\ref{EBAQJ}) holds.   
  Fix an arbitrary \(x \in A\). 
  
  If \(\theta(x) = 0\), then \(x \in E_w\) and taking \(y = x\) fulfills the requirement.
  
  If \(\theta(x) > 0\),  by   (\ref{EBAQJ}),     there exists \(x' \in E_w\) such that
  \[
  \|x - x'\| < d(x,E_w) + \theta(x) \le \tau\theta(x) + \theta(x) = (\tau+1)\theta(x).
  \]
  Set \(y = x'\),   then \(y \in A\) and \(\theta(y) = 0 \le \beta\theta(x)\) for any \(\beta \in (0,1)\). Thus (iv) holds with constant \(\alpha = \tau+1\) and an arbitrary \(\beta \in (0,1)\).
  
  
  
  	\noindent\textbf{(iv) $\Rightarrow$ (i)}.
  Assume there exist \(\alpha>0\) and \(\beta \in (0,1)\) satisfying  (\ref{EYZXJ}). Set \(\tau = \frac{\alpha}{1-\beta}\).
    Fix an arbitrary \(x \in A\). 
    
      If \(\theta(x)=0\), then \(x \in E_w\) and the conclusion is trivial. 
      
      Suppose \(\theta(x) > 0\). Starting from \(x_0 = x\), we recursively construct a sequence \(\{x_n\} \subseteq A\). 
  Assuming \(x_n\) is already obtained, if \(\theta(x_n) = 0\) we stop,  otherwise by condition (\ref{EYZXJ}) there exists \(x_{n+1} \in A\) such that
  \[
  \|x_n - x_{n+1}\| \le \alpha \theta(x_n), \qquad \theta(x_{n+1}) \le \beta \theta(x_n).
  \]
  Iterating these inequalities yields, for all steps,
  \[
  \theta(x_n) \le \beta^n \theta(x), \qquad \|x_n - x_{n+1}\| \le \alpha \beta^n \theta(x).
  \]
  Since \(0 < \beta < 1\), the series \(\sum_{n=0}^\infty \beta^n\) converges, hence \(\{x_n\}\) is a Cauchy sequence. As a closed subset of a Banach space, \(A\) is complete, so there exists a limit point \(\bar{x} \in A\) with \(x_n \to \bar{x}\). 
  By the Lipschitz continuity of \(\theta\),
  \[
  \theta(\bar{x}) = \lim_{n\to\infty} \theta(x_n) \le \lim_{n\to\infty} \beta^n \theta(x) = 0.
  \]
  Since \(\theta\) is nonnegative on \(A\), we obtain \(\theta(\bar{x}) = 0\), i.e., \(\bar{x} \in E_w\).   Then, we have
  \[ d(x,E_w) \le
  \|x - \bar{x}\| \le \sum_{n=0}^{\infty} \|x_n - x_{n+1}\|
  \le \sum_{n=0}^{\infty} \alpha \beta^n \theta(x)
  = \frac{\alpha}{1-\beta}\,\theta(x) = \tau \theta(x).
  \]
  This establishes (i).


	\noindent\textbf{(i) $\Rightarrow$ (v)}.
Suppose there exists \(\tau > 0\)  such that   (\ref{EBAQJ}) holds. Take arbitrary \(x \in A\) and \(e \in X\).

It follows from  \(x \in A\)  that \(d(x+e, A) \le \|x+e - x\| = \|e\|\).
  For any \(\varepsilon > 0\), choose \(a_\varepsilon \in A\) such that
\begin{equation}\label{eTHVDD1}
\|x+e - a_\varepsilon\| \le d(x+e, A) + \varepsilon \le \|e\| + \varepsilon,
\end{equation}
Since  \(\theta\) is Lipschitz continuous on \(X\) with constant \(L = \|F\|\),  by (\ref{eTHVDD1}), we have 
\[
|\theta(x+e) - \theta(a_\varepsilon)| \le L \|(x+e) - a_\varepsilon\| \le L(\|e\| + \varepsilon),
\]
hence
\begin{equation}\label{eTHVDD2}
	\theta(a_\varepsilon) \le \theta(x+e) + L(\|e\| + \varepsilon).
\end{equation}
 We conclude from  \(a_\varepsilon \in A\),  (\ref{EBAQJ})  and  (\ref{eTHVDD2})  that
\begin{equation}\label{eTHVDD3}
	d(a_\varepsilon, E_w) \le \tau \theta(a_\varepsilon) \le \tau\bigl(\theta(x+e) + L(\|e\| + \varepsilon)\bigr).
\end{equation}
  It follows from  (\ref{eTHVDD1})  and  (\ref{eTHVDD3})    that
\[
\begin{aligned}
	d(x+e, E_w) &\le \|x+e - a_\varepsilon\| + d(a_\varepsilon, E_w) \\
	&\le (\|e\| + \varepsilon) + \tau\bigl(\theta(x+e) + L(\|e\| + \varepsilon)\bigr) \\
	&= \tau \theta(x+e) + (1 + \tau L)(\|e\| + \varepsilon).
\end{aligned}
\]
Letting \(\varepsilon \to 0^+\), we obtain
\begin{equation}\label{eq:d-xe}
	d(x+e, E_w) \le \tau \theta(x+e) + (1 + \tau L)\|e\|.
\end{equation}
Since  \(\theta\) is Lipschitz continuous, \(\theta(x+e) \le \theta(x) + L\|e\|\). Substituting into \eqref{eq:d-xe} yields
\[
\begin{aligned}
	d(x+e, E_w) &\le \tau\bigl(\theta(x) + L\|e\|\bigr) + (1 + \tau L)\|e\| \\
	&= \tau \theta(x) + (2\tau L + 1)\|e\|.
\end{aligned}
\]
 Due to  \(\theta(x) \ge 0\) and \(\|e\| \ge 0\), we have
\[
\tau \theta(x) + (2\tau L + 1)\|e\| \le (\tau + 2\tau L + 1)(\theta(x) + \|e\|).
\]
Setting \(\kappa = \tau + 2\tau L + 1\), we obtain
\[
d(x+e, E_w) \le \kappa (\theta(x) + \|e\|),
\]
  which means that  (\ref{EBkkQJ}) holds.

	\noindent\textbf{(v) $\Rightarrow$ (i)}.
Suppose  that There exist \(\kappa > 0\) such that  (\ref{EBkkQJ}) holds.   
Take any \(x \in A\) and set \(e = 0\). From (\ref{EBkkQJ}) we immediately obtain
\[
d(x, E_w) = d(x+0, E_w) \le \kappa (\theta(x) + 0) = \kappa \theta(x).
\]
 
 
	\noindent\textbf{(i) $\Rightarrow$ (vi)}.
Assume  that   there exists $\tau>0$ such that  (\ref{EBAQJ}) holds.
Set $L=\|F\|$, and define
\[
\gamma := \max\{\,\tau,\;1+\tau L\,\}.
\]
Let  $x\in X$      be arbitrary.   For   any    $\varepsilon>0$,  there exists $a_\varepsilon\in A$ such that
 \begin{equation}\label{EvgyB1}	\|x-a_\varepsilon\| < d_A(x)+\varepsilon  .        
 \end{equation}
 It follows from    (\ref{EBAQJ})   and  (\ref{EvgyB1})    that   
 \begin{equation}\label{EvgyB2}
 	d(x,E_w) \le \|x-a_\varepsilon\| + d(a_\varepsilon,E_w) < d_A(x)+\varepsilon + \tau\,\theta(a_\varepsilon) .
 \end{equation}
Since  \(\theta\) is Lipschitz continuous on \(X\) with constant \(L = \|F\|\),  due to (\ref{EvgyB1}),  we have
\begin{equation} \label{EvgyB3}
	\theta(a_\varepsilon) \le \theta(x) + L\|a_\varepsilon-x\| 
	< \theta(x) + L\bigl(d_A(x)+\varepsilon\bigr).  
\end{equation}

Now we distinguish two cases according to the sign of $\theta(x)$.

\noindent\textbf{Case 1: $\theta(x)\ge 0$.}
Then $\max\{\theta(x),0\}=\theta(x)$. Substituting (\ref{EvgyB3}) into (\ref{EvgyB2}),
\[
\begin{aligned}
	d(x,E_w) &< d_A(x)+\varepsilon + \tau\bigl(\theta(x)+L(d_A(x)+\varepsilon)\bigr) \\
	&= \tau\theta(x) + (1+\tau L)d_A(x) + (1+\tau L)\varepsilon.
\end{aligned}
\]
Since $\tau\le\gamma$ and $1+\tau L\le\gamma$, and $\theta(x)\ge0,\;d_A(x)\ge0$, we have
\[
\tau\theta(x) + (1+\tau L)d_A(x) \le \gamma\theta(x) + \gamma d_A(x) = \gamma\bigl(\theta(x)+d_A(x)\bigr).
\]
Thus
\[
d(x,E_w) < \gamma\bigl(\theta(x)+d_A(x)\bigr) + (1+\tau L)\varepsilon.
\]

\noindent\textbf{Case 2: $\theta(x)<0$.}
Then $\max\{\theta(x),0\}=0$. Due to (\ref{EvgyB3})  and   $\theta(x)<0$, we have
$\theta(a_\varepsilon) < L\bigl(d_A(x)+\varepsilon\bigr)$.
Substituting into (\ref{EvgyB2}), we have
\[
\begin{aligned}
	d(x,E_w) &< d_A(x)+\varepsilon + \tau \cdot L\bigl(d_A(x)+\varepsilon\bigr) \\
	&= (1+\tau L)d_A(x) + (1+\tau L)\varepsilon.
\end{aligned}
\]
Since $1+\tau L\le\gamma$ and in this case $\max\{\theta(x),0\}+d_A(x)=d_A(x)$, we obtain
\[
d(x,E_w) < \gamma\bigl(\max\{\theta(x),0\}+d_A(x)\bigr) + (1+\tau L)\varepsilon.
\]

Combining both cases, for any $\varepsilon>0$ we have
\[
d(x,E_w) < \gamma\bigl(\max\{\theta(x),0\}+d_A(x)\bigr) + (1+\tau L)\varepsilon.
\]
Letting $\varepsilon\to0^+$, we get
\[
d(x,E_w) \le \gamma\bigl(\max\{\theta(x),0\}+d_A(x)\bigr).
\]
 This means that (\ref{EBkXJ}) is true.

  
 
 	\noindent\textbf{(vi) $\Rightarrow$ (i)}.
 Assume that  there exists $\gamma>0$ such that (\ref{EBkXJ}) is satisfied.
  Take any $x\in A$. Then $d_A(x)=0$  and   $\theta(x)\ge0$, and so $\max\{\theta(x),0\}=\theta(x)$. Substituting into (\ref{EBkXJ}) immediately yields
 \[
 d(x,E_w)\le \gamma\,\theta(x),\quad \forall x\in A.
 \]
  
 
 
  	\noindent\textbf{(i) $\Rightarrow$ (vii)}.
  Suppose there exists \(\tau > 0\) such that   (\ref{EBAQJ}) is true.
 Take an arbitrary \(x \in X\). For any \(a \in A\),   it follows from  (\ref{EBAQJ})  that
 \[
 d(x, E_w) \le \|x - a\| + d(a, E_w) \le \|x - a\| + \tau \, \theta(a).
 \]
 Set \(\kappa := \max\{\tau, 1\}\). Then  
$  \|x - a\| + \tau \, \theta(a) \le \kappa \bigl( \|x - a\| + \theta(a) \bigr),$  and so
   $$d(x, E_w) \le \kappa (\|x - a\| + \theta(a)),  \quad  \forall a \in A.  $$
 Taking the infimum over \(a \in A\) yields  
 \[
 d(x, E_w) \le \kappa \inf_{a \in A} \bigl( \theta(a) + \|x - a\| \bigr) = \kappa \, \zeta(x).
 \]
 Since \(x \in X\) is arbitrary, this proves (vii).
 
 
  	\noindent\textbf{(vii) $\Rightarrow$ (i)}.
 Suppose there exists \(\kappa > 0\) such that 
 \begin{equation}\label{ErtB1}
 	d(x, E_w) \le \kappa \, \zeta(x), \quad  \forall x \in X.
 \end{equation}
 Take an arbitrary \(x \in A\). From the definition of \(\zeta\) we have  
 \[
\zeta(x) = \inf_{a \in A} \bigl( \theta(a) + \|x - a\| \bigr) \le \theta(x) + \|x - x\| = \theta(x).
 \]
 Substituting this into    (\ref{ErtB1})   gives  
 \[
 d(x, E_w) \le \kappa \, \theta(x), \quad \forall x \in A.
 \]
  
	

	 	
	  
	 		\noindent\textbf{(i) $\Rightarrow$ (ix)}.
	 			 		Assume that there exists \(\tau>0\) such that   (\ref{EBAQJ}) is true.
	 	 	 		Take an arbitrary \(x\in A\setminus E_w\). Then \(\theta(x)>0\)  and	$	d:=d(x,E_w)>0 .$
	 		 For any \(\varepsilon>0\), there exists \(y_\varepsilon\in E_w\) satisfying
	 		\begin{equation}\label{Eapprox1}
	 			\|x-y_\varepsilon\|<d+\varepsilon .
	 		\end{equation}
	 		Since \(y_\varepsilon\in E_w\), we have \(\theta(y_\varepsilon)=0<\theta(x)\).  so \(y_\varepsilon\) belongs to the set over which the supremum defining \(m(x)\) is taken.   It follows from  (\ref{Eapprox1})   and    \(\theta(x)>0\)  that
	 		\begin{equation}\label{Eapprox2}
	 		\frac{\theta(x)-\theta(y_\varepsilon)}{\|x-y_\varepsilon\|}=\frac{\theta(x)}{\|x-y_\varepsilon\|}>\frac{\theta(x)}{d+\varepsilon}.
	 		\end{equation}
	 		 From (\ref{EBAQJ}) we obtain \(\theta(x)\ge d/\tau\), and so
	 	\begin{equation}\label{Eapprox3}
	 		\frac{\theta(x)}{d+\varepsilon}\ge\frac{d/\tau}{d+\varepsilon}=\frac{1}{\tau}\cdot\frac{d}{d+\varepsilon}.
	 		\end{equation}
	 		Therefore, for this \(y_\varepsilon\),  we conclude from  (\ref{Eapprox2})  and  (\ref{Eapprox3})  that
	 		\[
	 		m(x)\ge\frac{\theta(x)-\theta(y_\varepsilon)}{\|x-y_\varepsilon\|}>\frac{1}{\tau}\cdot\frac{d}{d+\varepsilon}.
	 		\]
	 		Letting \(\varepsilon\to0^{+}\), the right-hand side tends to \(\gamma:=\frac{1}{\tau}\),  hence \(m(x)\ge \gamma\). Since \(x\in A\setminus E_w\) was arbitrary, this proves (ix).
	 	 
	 		
	 			\noindent\textbf{(ix) $\Rightarrow$ (viii)}.
	 	Suppose   that there exists   \(\gamma >0\) such that
	 	\[
	 	m(x)\ge \gamma,\quad \forall x\in A\setminus E_w.
	 	\]
	 	Set \(c :=\dfrac{\gamma}{2}>0\). For any \(x\in A\setminus E_w\), since \(m(x)\ge \gamma > c\), by the definition of supremum, there exists \(y\in A\) satisfying \(\theta(y)<\theta(x)\) and
	 	\[
	 	\frac{\theta(x)-\theta(y)}{\|x-y\|} > c.
	 	\]
	 	Rearranging yields \(\theta(y) < \theta(x) - c \|x-y\|\).   It follows from \(\theta(y)<\theta(x)\)  that    \(y\neq x\). Therefore (viii) holds (with constant \(c =\dfrac{\gamma}{2}\)).
	 	
	\noindent\textbf{(viii) $\Rightarrow$ (i)}.
	 		Assume that  there exists   \(c>0\) such that     (\ref{EevL})   holds.
	 			We will prove that for all \(x\in A\),
	 		\[
	 		d(x,E_w)\le \frac{1}{c}\,\theta(x).
	 		\]
	 		Suppose, to the contrary, that there exists \(x_0\in A\) such that  
	 		\begin{equation}\label{Econtra1}
	 			d(x_0,E_w)>\frac{1}{c}\,\theta(x_0).
	 		\end{equation}
	 	 This yields  that	 \(\theta(x_0)>0\). 
	 		Consider the restriction \(f:=\theta|_A\colon A\to[0,\infty)\) on the complete metric space \(A\) (a closed subset of a Banach space). It is clear that \(f\) is continuous  and \(\inf_{x \in A} f(x)=0\).  
	 		Set
	 		\[
	 		\varepsilon:=f(x_0)=\theta(x_0)>0,\qquad \lambda:=\frac{\varepsilon}{c}=\frac{\theta(x_0)}{c}>0 .
	 		\]
	 		Then \(f(x_0)\le \inf_{x \in A} f(x) +\varepsilon\).    Lemma  \ref{ELdaz} (Ekeland's variational principle)   yields that  there exists \(\bar x\in A\) satisfying:
	 		\begin{itemize}
	 			\item[(a)] \(f(\bar x)\le f(x_0)\);
	 			\item[(b)] \(\|\bar x-x_0\|\le \lambda\);
	 			\item[(c)] for all \(x\in A\setminus\{\bar x\}\), \(f(x) > f(\bar x)-\dfrac{\varepsilon}{\lambda}\,\|x-\bar x\|\).
	 		\end{itemize}
	 		If \(\bar x \in  E_w\), then  it follows from  (b) that
	 		\[
	 		d(x_0,E_w)\le \|x_0-\bar x\|\le \lambda = \frac{1}{c}\theta(x_0),
	 		\]
	 	which contradicts (\ref{Econtra1}). Hence   \(\bar x \notin  E_w\),  and so   \(\theta(\bar x)=f(\bar x)>0\).   (c) can be rewritten as
	 		\begin{equation}\label{Econtra2}
	 			\forall x\in A\setminus\{\bar x\},\qquad \theta(x) > \theta(\bar x) - c\,\|x-\bar x\|,
	 		\end{equation}
	 		because \(\varepsilon/\lambda = c\).
	 			 		Now apply  (\ref{EevL})   to \(\bar x\in A\setminus E_w\),  there exists \(y_0 \in A\) with \(y_0 \neq \bar x\) such that
	 		\begin{equation}\label{Econtra3}
	 			\theta(y_0 ) \le \theta(\bar x) - c\,\|\bar x-y_0\|.
	 		\end{equation}
	 It follows from  (\ref{Econtra2})  that		$\theta(y_0) > \theta(\bar x) - c\,\|y_0-\bar x\|,$
	 		which contradicts (\ref{Econtra3}). 
	 Therefore
	 		\[
	 		d(x,E_w)\le \frac{1}{c}\,\theta(x), \quad \forall x\in A,
	 		\]
	 		which is exactly statement (i) with \(\tau = 1/c\).
	 		
	\noindent\textbf{(i) $\Rightarrow$ (x)}.
Assume that  there exists    \(\tau > 0\) such that   (\ref{EBAQJ})  is satisfied.    Take any $x\in A$. 
For any $y\in E_w$,  we have
\[
\|F(x)-F(y)\|=\|F(x-y)\|\le\|F\|\,\|x-y\|.
\]
Taking the infimum over all $y\in E_w$ on both sides yields
\[
\inf_{y\in E_w}\|F(x)-F(y)\|\le\|F\|\inf_{y\in E_w}\|x-y\|.
\]
This means that
\begin{equation}\label{Efty1}
	d \bigl(F(x),F(E_w)\bigr)\le\|F\|\,d(x,E_w).  
\end{equation}
It is easy to see that $F(E_w)=\operatorname{WMin}(F(A),C)$.
Combining this with   (\ref{EBAQJ})  and  (\ref{Efty1}), we get
\[
d \bigl(F(x),\operatorname{WMin}(F(A),C)\bigr)\le\|F\|\,d(x,E_w)   \le\|F\|\,\tau\,\theta(x).
\]
Set $c:=\|F\|\,\tau >0$. Since $x\in A$ was arbitrary, (x) follows.	 		
	 		
	 		
	 	\noindent\textbf{(iii) $\Rightarrow$ (xi)}.
	 If (iii) holds, then for all $x\in A\setminus E_w$,
	 \[
	 \frac{d_{\widehat{A}}(x)}{d(x,E_w)}\ge\frac{1}{\tau}.
	 \]
	 As $x\xrightarrow{A} E_w$, the limit inferior is at least $1/\tau>0$, so (xi) holds.		
	 		
		 	\noindent\textbf{(i) $\Rightarrow$ (xii)}.
Suppose that  there exists \(\tau > 0\) such that   (\ref{EBAQJ})  is satisfied.

Take arbitrary \(t \ge 0\) and \(x \in \Phi(t)\). Since \(x \in A\) and \(\theta(x) \le t\),   (\ref{EBAQJ})  yields
\[
d(x, E_w) \le \tau \theta(x) \le \tau t.
\]
Because \(E_w \subseteq \Phi(t)\), clearly \(\sup_{y \in E_w} d(y, \Phi(t)) = 0\). Hence the Hausdorff distance satisfies
\[
d_H(\Phi(t), E_w) = \sup_{x \in \Phi(t)} d(x, E_w) \le \tau t, \quad \forall t \ge 0.
\]
Taking \(t_0 = 1\) and \(\beta = \tau\), we have for all \(t \in [0, t_0]\) that \(d_H(\Phi(t), E_w) \le \beta t\). 	 		
	 		
	
	\noindent\textbf{(i) $\Rightarrow$ (xiii)}.
	By definition of $\Phi$, clearly $\Phi(0)=\theta^{-1}(0)\cap A=E_w$.
	For a fixed $x\in A$,
	\[
	\Phi^{-1}(x)=\{t\ge0:\theta(x)\le t\}=[\theta(x),+\infty),
	\]
and	so $d(0,\Phi^{-1}(x))=\inf\{|t-0|:t\ge\theta(x)\}=\theta(x)$.
	Therefore,      (\ref{EBkmJ})    is equivalent to
\begin{equation}\label{EViiiD}
	d(x,E_w)\le\kappa\,\theta(x),\quad\forall x\in U\cap A.
 \end{equation}
	If (i) holds, taking $U=X$ shows that (xiii) holds.
	
	
	
	
	
	
	
	
	
	
 

		
		
	\noindent\textbf{(x) $\Rightarrow$ (i)}.  Suppose that   there exist $\mu>0$ such that $\|F(x)\|\ge\mu\|x\|$ for any $x\in A-A$.    Fix an arbitrary $x\in A$.
					 By linearity of $F$, for any $y\in E_w    \subseteq   A$,  we have
			  \[
			 \|F(x)-F(y)\|=\|F(x-y)\|\ge\mu\|x-y\|,
			 \]
 and so 
		\[
		\|x-y\|\le\frac{1}{\mu}\,\|F(x)-F(y)\|.
		\]
		Taking the infimum over all $y\in E_w$, we get
		\begin{equation}\label{EfghB1}
		d(x,E_w)=\inf_{y\in E_w}\|x-y\|
		\le\inf_{y\in E_w}\frac{1}{\mu}\,\|F(x)-F(y)\|
		=\frac{1}{\mu}\,\inf_{y\in E_w}\|F(x)-F(y)\|.
	 \end{equation}
The expression $\inf_{y\in E_w}\|F(x)-F(y)\|$ is exactly $d \bigl(F(x),F(E_w)\bigr)$.  By  $F(E_w)=\operatorname{WMin}(F(A),C)$  and    (\ref{EfghB1}),  we get  
		\begin{equation}\label{EfghB2}
		d(x,E_w)\le\frac{1}{\mu}\,d \bigl(F(x),\operatorname{WMin}(F(A),C)\bigr).  
		 \end{equation}
		  (x) provides a constant $c>0$   such that    
  $$ d \bigl(F(x),\operatorname{WMin}(F(A),C)\bigr)     \le c \; \theta(x),   \quad  \forall x\in A.$$ 
	This together with   (\ref{EfghB2})    implies that  
	  we get
		\[
		d(x,E_w)\le\frac{1}{\mu}\cdot  c  \,\theta(x)=\frac{c}{\mu}\,\theta(x).
		\]
		Set $\tau:=c/\mu>0$,  then
		\[
		d(x,E_w)\le\tau\,\theta(x).
		\]
		This holds for all $x\in A$, which is exactly statement (i).
 
	
		 	\noindent\textbf{(xii) $\Rightarrow$ (i)}.
		Suppose there exist \(t_0 > 0\) and \(\beta > 0\) such that 
		\begin{equation}\label{EthhB1}
			d_H(\Phi(t), E_w) \le \beta t,  \quad  \forall t \in [0, t_0].
		\end{equation}
		
 
	Take any \(x \in A\) and set \(t = \theta(x)\). We distinguish two cases:
	
  If \(t \le t_0\), then \(x \in \Phi(t)\). By   (\ref{EthhB1}),  we have 
		\[
		d(x, E_w) \le \sup_{z \in \Phi(t)} d(z, E_w) = d_H(\Phi(t), E_w) \le \beta t = \beta \theta(x).
		\]
		
		 If \(t > t_0\), since \(A\) is bounded, let \(\eta = \operatorname{diam}(A) = \sup_{a,b \in A} \|a - b\| < +\infty\). Because \(x \in A\) and \(E_w \subseteq A\), we have \(d(x, E_w) \le \eta\). Combining this with \(t > t_0\) yields
		\[
		d(x, E_w) \le \eta < \frac{\eta}{t_0} t = \frac{\eta}{t_0} \theta(x).
		\]
 
	Setting \(\tau = \max\{\beta, \eta/t_0\}\), we obtain \(d(x, E_w) \le \tau \theta(x)\) for all \(x \in A\). 
	
	\noindent\textbf{(xi) $\Rightarrow$ (iii)}.
	From (xi), set
	\[
	\gamma := \liminf_{\substack{x\xrightarrow{A}E_w\\ x\notin E_w}} \frac{d_{\widehat{A}}(x)}{d(x,E_w)} > 0.
	\]
	Hence there exists $\eta>0$ such that for all $x\in A$ with $0<d(x,E_w)<\eta$,
	\[
	\frac{d_{\widehat{A}}(x)}{d(x,E_w)} \ge \frac{\gamma}{2} .
	\]
	Equivalently,
\begin{equation}\label{HYTz1}
	d(x,E_w) \le \frac{2}{\gamma}\, d_{\widehat{A}}(x), \qquad \forall\,x\in A,\; 0<d(x,E_w)<\eta.  
 \end{equation}
		Let $Q := \{x\in A : d(x,E_w)\ge \eta\}$. Since $A$ is compact and the distance function is continuous, $Q$ is compact.
	
	We now show that $\inf_{x\in Q} d_{\widehat{A}}(x) > 0$.
	Suppose that  $\inf_{x\in Q} d_{\widehat{A}}(x) = 0$.   By the compactness of $Q$ and continuity of $d_{\widehat{A}}$,  there exists  $\bar{x}\in Q \subseteq A$ such that
	$$d_{\widehat{A}}(\bar{x})=  \inf_{x\in Q} d_{\widehat{A}}(x) = 0 .$$
	This together with  the  closedness of  $\widehat{A}$  implies that $\bar{x}\in \widehat{A}$.
	 Moreover $\bar{x}\in A$, so $\bar{x}\in A\cap\widehat{A}=E_w$. This contradicts $d(\bar{x},E_w)\ge \eta>0$. Hence
\begin{equation}\label{HYTz2}
	\alpha := \inf_{x\in Q} d_{\widehat{A}}(x) >0.  
 \end{equation}
	
	Take any $x\in A$, and consider two cases:
 
 Case 1:  $d(x,E_w)<\eta$.  It follows from  (\ref{HYTz1})  that
		\[
		d(x,E_w)\le \frac{2}{\gamma} d_{\widehat{A}}(x).
		\]
		
	  Case 2:   $d(x,E_w)\ge \eta$, i.e., $x\in Q$. Using  (\ref{HYTz2})    and the diameter $\beta :=\operatorname{diam}(A)=\sup_{u,v\in A}\|u-v\|<\infty$,  we have 
		\[
		d(x,E_w) \le \beta \le \frac{\beta}{\alpha}\, d_{\widehat{A}}(x).
		\]
 	Taking $\tau = \max\left\{\dfrac{2}{\gamma},\; \dfrac{\beta}{\alpha}\right\}$, we obtain for all $x\in A$,
	\[
	d(x,E_w) \le \tau\, d_{\widehat{A}}(x).
	\]
	This is the global error bound (iii).
	
	\noindent\textbf{(xiii) $\Rightarrow$ (i)}.
	Assume $A$ is compact   and there exist $\kappa>0$ and an open set $U\supseteq E_w$ such that
  (\ref{EBkmJ})  is satisfied.   
	Set $H:=A\setminus U$.   Since   $A$ is compact, $U$ is open and $E_w \subseteq U$,  we get that  $H$ is compact   and $H\cap E_w= \emptyset$.
	The function $\theta$ is continuous on the compact set $H$, hence attains its minimum $m:=\min_{x\in H}\theta(x)$.
	If $m=0$, then there exists $\bar x\in H$ with $\theta(\bar x)=0$, implying $\bar x\in E_w$, contradicting $H\cap E_w= \emptyset$. Therefore $m>0$.
	
	For any $x\in H$, since $A$ is bounded, $\beta =\operatorname{diam}(A)<\infty$, and so $d(x,E_w)\le \beta$. Together with $\theta(x)\ge m$, we have 
	\begin{equation} \label{EHYZ2}
		d(x,E_w)\le \frac{\beta}{m}\,\theta(x).  
	\end{equation}
	
	Take any $x\in A$.  If $x\in U$, noting that  (\ref{EBkmJ})  is  equivalent to (\ref{EViiiD}), we get  
	   $d(x,E_w)\le\kappa\,\theta(x)$.
  
 If $x\in H$,  it follows from  (\ref{EHYZ2}) that $d(x,E_w)\le (\beta/m)\,\theta(x)$.
 	Taking $\tau:=\max\{\kappa,\; \beta/m\}$, the global error bound holds:
	\[
	d(x,E_w)\le\tau\,\theta(x),\qquad \forall x\in A.
	\]
	This completes the proof. 
	\end{proof}


\begin{thebibliography}{99}
\bibitem[AE1984]{AE1984} Aubin, J.P., Ekeland, I.: Applied Nonlinear Analysis. Pure and Applied Mathematics, New York (1984)
\bibitem[DF2024]{DF2024} Durea, M., Florea, E.A.: Cone-compactness of a set and applications to set-equilibrium problems. J. Optim. Theory Appl. \textbf{200}, 1286–1308 (2024)
\bibitem[Han2022]{Han2022} Han, Y.: Some characterizations of a nonlinear scalarizing function via oriented distance function. Optimization \textbf{71}, 4785–4817 (2022)
\bibitem[Han2025]{Han2025} Han, Y.: Cone sequential compactness of a set and an application to set optimization problems. J. Optim. Theory Appl. \textbf{206}, 59 (2025)
\bibitem[Han2026]{Han2026} Han, Y.: Nonlinear scalarization for Lagrangian duality in set optimization: saddle point characterizations. Preprint (2026). DOI: 10.13140/RG.2.2.30758.69440
\bibitem[Hiriart1]{Hiriart1} Hiriart-Urruty, J.B.: New concepts in nondifferentiable programming. Bull. Soc. Math. France Mém. \textbf{60}, 57–85 (1979)
\bibitem[Hiriart2]{Hiriart2} Hiriart-Urruty, J.B.: Tangent cones, generalized gradients and mathematical programming in Banach spaces. Math. Oper. Res. \textbf{4}, 79–97 (1979)
\bibitem[KTZ2015]{KTZ2015} Khan, A.A., Tammer, C., Zălinescu, C.: Set-valued Optimization: An Introduction with Applications. Springer, Heidelberg (2015)
\bibitem[LiuNgYang2009]{LiuNgYang2009} Liu, C.G., Ng, K.F., Yang, W.H.: Merit functions in vector optimization. Math. Program. Ser. A \textbf{119}, 215–237 (2009)
\bibitem[Luc]{Luc} Luc, D.T.: Theory of Vector Optimization. Lecture Notes in Economics and Mathematical Systems, vol. 319, Springer, Berlin (1989)
\bibitem[XuLi2016]{XuLi2016} Xu, Y.D., Li, S.J.: A new nonlinear scalarization function and applications. Optimization \textbf{65}, 207–231 (2016)
\bibitem[Zaffaroni2003]{Zaffaroni2003} Zaffaroni, A.: Degrees of efficiency and degrees of minimality. SIAM J. Control Optim. \textbf{42}, 1071–1086 (2003)
\bibitem[Zalinescu1978]{Zalinescu1978} Zălinescu, C.: A generalization of the Farkas lemma and applications to convex programming. J. Math. Anal. Appl. \textbf{66}, 651–678 (1978)
\bibitem[Zalinescu2002]{Zalinescu2002} Zălinescu, C.: Convex Analysis in General Vector Spaces. World Scientific, Singapore (2002)
\end{thebibliography}
\end{document}
